\section{Accounting as Conservation Laws}
\label{sec:conservation}

This section develops the mathematics. Everything here is stated so that
the Lean formalization of \cref{sec:proofs} is a direct transcription:
the definitions below have the same names as their Lean counterparts in
\texttt{Ratio.Core}.

\subsection{The conserved-dimension basis}

Fix a (possibly infinite, dynamically extensible) set $\rdim$ of
\emph{conserved dimensions}. A dimension is any quantity that must be
neither created nor destroyed by a transaction --- a currency's minor
units, a security position, a tax lot, a P\&L component. Accounting
state is an integer-valued vector over this basis.

\begin{definition}[State / net-effect vector]
\label{def:vec}
A state is a function $\rvec : \rdim \to \mathbb{Z}$, written
$\rvec \in \mathbb{Z}^{\rdim}$. Addition and negation are pointwise,
\[
  (\rvec + \mathbf{w})(d) \;=\; \rvec(d) + \mathbf{w}(d),
  \qquad
  (-\rvec)(d) \;=\; -\,\rvec(d),
\]
and $\rzero$ is the everywhere-zero vector $\rzero(d) = 0$.
\end{definition}

We work over $\mathbb{Z}$, never $\mathbb{R}$ or floating point. Money is
represented in integer minor units (e.g.\ cents): this is the only
representation under which conservation is \emph{exact} and replay is
\emph{deterministic} across architectures. Floating-point addition is
neither associative nor exactly invertible~\cite{goldberg1991floating},
so it cannot carry a conservation law at all. The group-theoretic view of
double-entry is due to Ellerman~\cite{ellerman1985mathematics}; we
generalize it from a single balancing dimension to an arbitrary basis.

\begin{proposition}[$\mathbb{Z}^{\rdim}$ is a commutative group]
\label{prop:group}
$(\mathbb{Z}^{\rdim}, +, \rzero, -)$ is an abelian group: $+$ is
associative and commutative, $\rzero$ is a two-sided identity, and $-\rvec$
is a two-sided inverse.
\leanref{Ratio.Vec.add\_comm, add\_assoc, add\_zero, add\_neg}
\end{proposition}

This is the algebraic substrate. It is unremarkable as mathematics ---
which is the point. The kernel does nothing but exploit the group
structure; all the difficulty of real accounting lives in choosing
$\rdim$ and interpreting the result, and that difficulty is pushed out of
the kernel entirely (\cref{sec:control}).

\subsection{Transactions and conservation}

\begin{definition}[Conservation / balance]
\label{def:conserves}
A vector $\rvec$ \emph{conserves value} --- equivalently, a transaction
with net effect $\rvec$ is \emph{balanced} --- iff
\[
  \rconserves(\rvec) \;:\Longleftrightarrow\; \rvec = \rzero,
  \qquad\text{i.e.}\qquad \forall d \in \rdim,\; \rvec(d) = 0 .
\]
\end{definition}

This single predicate \emph{is} the kernel's only invariant. A
transaction is a movement of value across dimensions; it is admissible
exactly when nothing leaks. \Cref{fig:vector} shows the picture: a
two-line cash purchase as a vector over a four-dimension basis, summing
to $\rzero$.

% A transaction as an integer vector over a conserved-dimension basis,
% netting to zero. Example: buy \$1{,}200 of equity, settling cash, with
% a brokerage fee booked to an expense dimension --- one balanced vector.
\begin{figure}[t]
\centering
\begin{tikzpicture}[font=\footnotesize, >={Stealth[length=2mm]}]
\matrix (m) [matrix of nodes,
    nodes={anchor=center, minimum height=6.5mm, minimum width=24mm,
           draw=black!25, inner sep=2pt, align=center},
    column sep=-\pgflinewidth, row sep=-\pgflinewidth,
    row 1/.style={nodes={fill=black!8, font=\bfseries}}]{
  dimension $d$ & net $\rvec(d)$ (minor units) \\
  Cash (USD)            & $-121500$ \\
  Equity position (USD) & $+120000$ \\
  Brokerage fee (USD)   & $+1500$ \\
};
% Sum bracket
\node[right=6mm of m, align=left] (sum) {%
  $\displaystyle\sum_{d\in\rdim}\rvec(d)=0$\\[2pt]
  {\scriptsize $\Rightarrow \rconserves(\rvec)$: the}\\[-1pt]
  {\scriptsize transaction is \emph{balanced}.}};
\draw[->] (m.east|-m-2-2) -- ++(0.5,0) |- (sum.west);
\draw[->] (m.east|-m-3-2) -- ++(0.5,0) |- (sum.west);
\draw[->] (m.east|-m-4-2) -- ++(0.5,0) |- (sum.west);
\end{tikzpicture}
\caption{A transaction is one integer vector $\rvec\in\mathbb{Z}^{\rdim}$
over a basis of conserved dimensions. Here a \$1{,}200 equity purchase
plus a \$15 fee settles against cash; the components net to zero across
the basis, so $\rconserves(\rvec)$ holds and the kernel admits it. Classical
double-entry is the one-dimensional shadow ($|\rdim|=1$): a list of signed
legs summing to zero.}
\label{fig:vector}
\end{figure}


\paragraph{Double-entry is the one-dimensional shadow.} Classical
double-entry bookkeeping is the special case $|\rdim| = 1$. A transaction
there is a list of signed split amounts $[a_1, \dots, a_n]$, ``balanced''
when $\sum_i a_i = 0$ --- the familiar ``debits equal credits''. The
canonical two-leg posting $[a, -a]$ balances because $a + (-a) = 0$, the
group inverse law. Ratio keeps double-entry as a derived, checked
instance (\cref{thm:double-entry}) but its kernel reasons over the full
$\mathbb{Z}^{\rdim}$, so a single FX trade that touches two currencies and
a realized-gain component is \emph{one} balanced transaction over three
dimensions rather than an awkward pile of bridging postings.

\begin{proposition}[Conserving vectors form a subgroup]
\label{prop:subgroup}
The set $\{\rvec : \rconserves(\rvec)\}$ contains $\rzero$ and is closed
under $+$ and $-$. Concretely:
$\rconserves(\rzero)$; and if $\rconserves(\mathbf{a})$ and
$\rconserves(\mathbf{b})$ then $\rconserves(\mathbf{a}+\mathbf{b})$ and
$\rconserves(-\mathbf{a})$.
\leanref{Ratio.conserves\_zero, conserves\_add, conserves\_neg}
\end{proposition}

\Cref{prop:subgroup} is the load-bearing fact. It says that balance is
\emph{compositional}: combine two balanced transactions and the result
is balanced; reverse a balanced transaction and it is still balanced.
Everything the ledger does is built from these two closures.

\subsection{The ledger as a monoidal fold}

\begin{definition}[Ledger]
\label{def:ledger}
Given an append-only log of transactions
$\textit{log} = [\rvec_1, \dots, \rvec_n]$, the ledger state is the
left fold
\[
  \rledger(\textit{log}) \;=\;
  \big(\cdots(\rzero + \rvec_1) + \rvec_2 \cdots\big) + \rvec_n
  \;=\; \sum_{i=1}^{n} \rvec_i .
\]
\end{definition}

Because $+$ is associative with identity $\rzero$
(\cref{prop:group}), $(\mathbb{Z}^{\rdim}, +, \rzero)$ is a commutative
monoid and the fold is well-defined independent of association --- it is
a \emph{monoidal aggregation}~\cite{hutton1999fold}. \Cref{fig:fold} renders this as the
commuting picture that makes the conservation law inevitable: folding the
log and checking balance commute, because every accumulator the fold
passes through is itself balanced.

% The ledger fold as a walk through the conserving subgroup. Each
% partial accumulator is balanced, so the final ledger state is balanced
% --- the inductive content of Theorem (conservation law).
\begin{figure}[t]
\centering
\begin{tikzpicture}[
    font=\footnotesize, >={Stealth[length=2mm]},
    acc/.style={draw, rounded corners, fill=green!10, minimum height=8mm,
                inner sep=3pt, align=center},
    txn/.style={font=\scriptsize},
]
\node[acc] (a0) {$\rzero$};
\node[acc, right=14mm of a0] (a1) {$\rzero{+}\rvec_1$};
\node[acc, right=14mm of a1] (a2) {$\cdots{+}\rvec_2$};
\node[right=11mm of a2] (dots) {$\cdots$};
\node[acc, right=10mm of dots] (an) {$\rledger(\textit{log})$};

\draw[->] (a0) -- node[above, txn] {$+\rvec_1$} (a1);
\draw[->] (a1) -- node[above, txn] {$+\rvec_2$} (a2);
\draw[->] (a2) -- node[above, txn] {$+\rvec_3$} (dots);
\draw[->] (dots) -- node[above, txn] {$+\rvec_n$} (an);

% Invariant annotation under every accumulator.
\node[font=\scriptsize\itshape, text=green!45!black, below=2.5mm of a0] {$\rconserves$};
\node[font=\scriptsize\itshape, text=green!45!black, below=2.5mm of a1] {$\rconserves$};
\node[font=\scriptsize\itshape, text=green!45!black, below=2.5mm of a2] {$\rconserves$};
\node[font=\scriptsize\itshape, text=green!45!black, below=2.5mm of an] {$\rconserves$};

\node[font=\scriptsize, align=center, below=12mm of a2]
  (note) {Every $\rvec_i$ conserves $\Rightarrow$ each step stays in the conserving subgroup
          (closure under $+$) $\Rightarrow$ by induction $\rconserves(\rledger(\textit{log}))$.};
\end{tikzpicture}
\caption{The ledger is a monoidal fold $\mathrm{foldl}\,(+)\,\rzero$ over
the append-only log. Conserving vectors form a subgroup
(\cref{prop:subgroup}); the seed $\rzero$ conserves, and adding a
conserving $\rvec_i$ keeps the accumulator conserving, so an induction on
the log proves the final state conserves (\cref{thm:conservation}).
Balance is preserved \emph{structurally} by the fold, not re-checked at
the end.}
\label{fig:fold}
\end{figure}


The central theorem, proved in \cref{sec:proofs}, is that this fold never
leaves the conserving subgroup:
\[
  \big(\forall i,\; \rconserves(\rvec_i)\big)
  \;\Longrightarrow\;
  \rconserves\!\big(\rledger(\textit{log})\big).
\]
In words: \emph{if every transaction balances, the whole ledger
balances} --- and it does so by a structural argument over the fold, not
by re-summing and hoping. This is what lets Ratio treat the ledger as
trustworthy infrastructure: internal consistency, replayability, and
auditability fall out of one closed-form invariant rather than out of
defensive checks scattered through application code.
